\documentclass[11pt,letterpaper]{article}
\usepackage[utf8]{inputenc}
\usepackage[margin=1in]{geometry}
\usepackage{tikz}
\usetikzlibrary{shapes.geometric, arrows.meta, positioning, calc, shadows}
\usepackage{booktabs}
\usepackage{longtable}
\usepackage{xcolor}
\usepackage{hyperref}
\usepackage{fancyhdr}
\usepackage{graphicx}

% Colors
\definecolor{routercolor}{RGB}{52, 73, 94}
\definecolor{tvcolor}{RGB}{155, 89, 182}
\definecolor{computercolor}{RGB}{41, 128, 185}
\definecolor{printercolor}{RGB}{39, 174, 96}
\definecolor{voipcolor}{RGB}{230, 126, 34}
\definecolor{iotcolor}{RGB}{241, 196, 15}
\definecolor{mobilecolor}{RGB}{231, 76, 60}
\definecolor{warningcolor}{RGB}{231, 76, 60}
\definecolor{goodcolor}{RGB}{39, 174, 96}
\definecolor{networkcolor}{RGB}{52, 152, 219}

\pagestyle{fancy}
\fancyhf{}
\rhead{Network Analysis Report}
\lhead{Superdooper 5G Network}
\rfoot{Page \thepage}

\title{\textbf{Network Analysis Report}\\[0.5em]\large Superdooper 5G (192.168.50.0/24)}
\author{Generated by Claude Code with WireMCP}
\date{\today}

\begin{document}

\maketitle
\tableofcontents
\newpage

%==============================================================================
\section{Executive Summary}
%==============================================================================

This report analyzes the \textbf{Superdooper 5G} WiFi network to identify:
\begin{itemize}
    \item Devices causing network dropouts
    \item Bandwidth-hogging devices
    \item Network topology and device inventory
    \item Optimization recommendations
\end{itemize}

\subsection{Key Findings}

\begin{itemize}
    \item \textcolor{goodcolor}{\textbf{Router Performance:}} ASUS RT-AX3000 operating normally (0\% packet loss, 2.3ms avg latency)
    \item \textcolor{goodcolor}{\textbf{WiFi Signal:}} Good quality (-63 dBm RSSI, 25 dB SNR, 648 Mbps)
    \item \textcolor{warningcolor}{\textbf{Potential Issue:}} Tenda device at .11 is offline/unresponsive - may be a dead range extender
    \item \textcolor{goodcolor}{\textbf{Protocol Mix:}} 97\% TCP/TLS (encrypted), minimal broadcast traffic
    \item \textbf{Total Devices:} 14 active devices discovered
\end{itemize}

%==============================================================================
\section{Network Topology}
%==============================================================================

\begin{center}
\begin{tikzpicture}[
    router/.style={rectangle, rounded corners, draw=routercolor, fill=routercolor!20,
                   thick, minimum width=3cm, minimum height=1.2cm, drop shadow},
    device/.style={rectangle, rounded corners, draw=#1, fill=#1!15,
                   thick, minimum width=2.5cm, minimum height=0.9cm, font=\small},
    extender/.style={rectangle, rounded corners, draw=warningcolor, fill=warningcolor!15,
                     thick, dashed, minimum width=2.2cm, minimum height=0.8cm, font=\small},
    connection/.style={->, >=Stealth, thick, networkcolor},
    label/.style={font=\tiny, text=gray}
]

% Router (center)
\node[router] (router) at (0,0) {\textbf{ASUS RT-AX3000}\\192.168.50.1};

% Internet cloud
\node[draw, cloud, cloud puffs=10, cloud puff arc=120, aspect=2,
      minimum width=2.5cm, minimum height=1.2cm, fill=blue!10]
      (internet) at (0,3) {Internet};
\draw[connection, very thick] (router) -- (internet);

% TVs (left side)
\node[device=tvcolor] (rokutv) at (-5,-2) {Roku TV 70"\\192.168.50.52};
\node[device=tvcolor] (googletv) at (-5,-3.5) {Google TV\\192.168.50.151};

% Computers (right side)
\node[device=computercolor] (macbook) at (5,-2) {MacBook Pro\\192.168.50.231};
\node[device=computercolor] (yoga) at (5,-3.5) {Lenovo Yoga\\192.168.50.37};
\node[device=computercolor] (hp) at (5,-5) {HP Roaster6\\192.168.50.158};

% IoT/Smart devices (bottom left)
\node[device=iotcolor] (amazon) at (-4,-5.5) {Amazon Echo\\192.168.50.40};
\node[device=iotcolor] (soundbar) at (-4,-7) {Yamaha ATS-1090\\192.168.50.195};

% VoIP (bottom center)
\node[device=voipcolor] (voip1) at (0,-5.5) {MagicJack\\192.168.50.124};
\node[device=voipcolor] (voip2) at (0,-7) {MagicJack\\192.168.50.150};
\node[device=voipcolor] (voip3) at (0,-8.5) {MagicJack\\192.168.50.201};

% Printer (bottom right)
\node[device=printercolor] (printer) at (4,-7) {Brother Printer\\192.168.50.114};

% Tenda (problematic - dashed)
\node[extender] (tenda) at (-2,-2) {Tenda (OFFLINE)\\192.168.50.11};

% Mobile/Unknown
\node[device=mobilecolor] (mobile) at (2,-3.5) {Mobile Device\\192.168.50.92};

% Connections
\draw[connection] (router) -- (rokutv);
\draw[connection] (router) -- (googletv);
\draw[connection] (router) -- (macbook);
\draw[connection] (router) -- (yoga);
\draw[connection] (router) -- (hp);
\draw[connection] (router) -- (amazon);
\draw[connection] (router) -- (soundbar);
\draw[connection] (router) -- (voip1);
\draw[connection] (router) -- (voip2);
\draw[connection] (router) -- (voip3);
\draw[connection] (router) -- (printer);
\draw[connection] (router) -- (mobile);
\draw[connection, dashed, warningcolor] (router) -- (tenda);

% Legend
\node[font=\footnotesize] at (0,-10) {
\begin{tabular}{cccccc}
\tikz\draw[fill=tvcolor!20, draw=tvcolor] (0,0) rectangle (0.3,0.2); TVs &
\tikz\draw[fill=computercolor!20, draw=computercolor] (0,0) rectangle (0.3,0.2); Computers &
\tikz\draw[fill=iotcolor!20, draw=iotcolor] (0,0) rectangle (0.3,0.2); IoT &
\tikz\draw[fill=voipcolor!20, draw=voipcolor] (0,0) rectangle (0.3,0.2); VoIP &
\tikz\draw[fill=printercolor!20, draw=printercolor] (0,0) rectangle (0.3,0.2); Printer &
\tikz\draw[fill=warningcolor!20, draw=warningcolor, dashed] (0,0) rectangle (0.3,0.2); Problem
\end{tabular}
};

\end{tikzpicture}
\end{center}

%==============================================================================
\section{Device Inventory}
%==============================================================================

\begin{longtable}{@{}llllp{3.5cm}@{}}
\toprule
\textbf{IP Address} & \textbf{MAC Address} & \textbf{Vendor} & \textbf{Type} & \textbf{Device Name} \\
\midrule
\endhead

192.168.50.1 & f0:2f:74:6e:5b:f8 & ASUS & Router & RT-AX3000-5BF8 \\
192.168.50.11 & c8:3a:35:3a:c2:35 & Tenda & Extender? & \textcolor{warningcolor}{OFFLINE} \\
192.168.50.37 & 14:ac:60:33:f6:51 & Lenovo & Laptop & TG-Yoga \\
192.168.50.40 & 48:b4:23:35:30:25 & Amazon & IoT & Echo/Fire device \\
192.168.50.52 & f0:a3:b2:b3:14:fd & Roku & Smart TV & 70" onn Roku TV \\
192.168.50.92 & 16:cb:af:68:84:dd & Private & Mobile & Unknown device \\
192.168.50.114 & 00:80:92:d6:6a:99 & Brother & Printer & MFC-L8850CDW \\
192.168.50.124 & 6c:33:a9:70:8d:75 & MagicJack & VoIP & Phone 1 \\
192.168.50.150 & 6c:33:a9:7b:6b:5e & MagicJack & VoIP & Phone 2 \\
192.168.50.151 & 02:58:a5:5a:ac:2a & Private & Smart TV & GoogleTV9356 \\
192.168.50.158 & dc:4a:3e:79:bb:7f & HP & Computer & Roaster6 \\
192.168.50.195 & 00:22:6c:18:e2:72 & Yamaha & Audio & ATS-1090 Soundbar \\
192.168.50.201 & 6c:33:a9:3e:20:48 & MagicJack & VoIP & Phone 3 \\
192.168.50.231 & 06:fe:91:16:fb:2d & Apple & Computer & MacBook Pro \\

\bottomrule
\end{longtable}

%==============================================================================
\section{Traffic Analysis}
%==============================================================================

\subsection{Protocol Distribution}

During the 20-second capture period:

\begin{center}
\begin{tabular}{lrr}
\toprule
\textbf{Protocol} & \textbf{Frames} & \textbf{Percentage} \\
\midrule
TCP/TLS (Encrypted) & 1,348 & 97.7\% \\
UDP (General) & 26 & 1.9\% \\
mDNS (Service Discovery) & 6 & 0.4\% \\
ARP & 2 & 0.1\% \\
\bottomrule
\end{tabular}
\end{center}

\textcolor{goodcolor}{\textbf{Good:}} Traffic is predominantly encrypted TCP/TLS with minimal broadcast traffic.

\subsection{Top Bandwidth Consumers}

30-second capture results:

\begin{center}
\begin{tabular}{llrp{4cm}}
\toprule
\textbf{Remote IP} & \textbf{Service} & \textbf{Data} & \textbf{Notes} \\
\midrule
34.36.57.103 & Google Cloud & 620 KB & Claude API traffic \\
160.79.104.10 & Unknown & 71 KB & General traffic \\
34.230.187.191 & AWS & 35 KB & Cloud services \\
142.250.191.10 & Google & 7.7 KB & Google services \\
\bottomrule
\end{tabular}
\end{center}

\subsection{Local Device Activity}

\begin{center}
\begin{tabular}{llp{6cm}}
\toprule
\textbf{Device} & \textbf{Status} & \textbf{Notes} \\
\midrule
Roku TV (.52) & Idle & Currently on HDMI 3 (Fire Stick) \\
Google TV (.151) & Idle & 16 bytes/sec (mDNS only) \\
MagicJack x3 & Idle & No active calls, no VoIP traffic \\
Soundbar (.195) & Idle & Periodic mDNS announcements \\
\bottomrule
\end{tabular}
\end{center}

%==============================================================================
\section{WiFi Signal Quality}
%==============================================================================

\begin{center}
\begin{tabular}{lll}
\toprule
\textbf{Metric} & \textbf{Value} & \textbf{Status} \\
\midrule
RSSI (Signal Strength) & -63 dBm & \textcolor{goodcolor}{Good} (threshold: $>$-70) \\
Noise Level & -88 dBm & \textcolor{goodcolor}{Good} (threshold: $<$-85) \\
SNR (Signal-to-Noise) & 25 dB & \textcolor{goodcolor}{Good} (threshold: $>$20) \\
Transmit Rate & 648 Mbps & \textcolor{goodcolor}{Excellent} \\
MCS Index & 6 & Moderate (max: 11) \\
PHY Mode & 802.11ax & WiFi 6 \\
Channel & 36 (5GHz, 80MHz) & Some congestion \\
\bottomrule
\end{tabular}
\end{center}

%==============================================================================
\section{Related Networks}
%==============================================================================

The following networks are saved on this device:

\begin{center}
\begin{tabular}{llp{5cm}}
\toprule
\textbf{Network Name} & \textbf{Band} & \textbf{Notes} \\
\midrule
Superdooper 5G & 5 GHz & Current network (WPA2/WPA3) \\
Test & Unknown & Not currently broadcasting \\
SuperdooperG & Unknown & Not in saved networks \\
Tenda\_6B0465\_5G & 5 GHz & Possibly the Tenda device at .11 \\
\bottomrule
\end{tabular}
\end{center}

\textbf{Note:} "Test" and "SuperdooperG" networks were not detected during the scan. They may be:
\begin{itemize}
    \item Disabled on the router
    \item Hidden SSIDs
    \item 2.4 GHz bands of the same router (common naming convention)
\end{itemize}

%==============================================================================
\section{Potential Dropout Causes}
%==============================================================================

Based on the analysis, here are the potential causes of network dropouts:

\subsection{High Priority Issues}

\begin{enumerate}
    \item \textcolor{warningcolor}{\textbf{Tenda Device at 192.168.50.11 (OFFLINE)}}
    \begin{itemize}
        \item Device is in ARP table but completely unresponsive (100\% packet loss)
        \item Likely a dead/crashed WiFi range extender
        \item If it intermittently comes online, it could cause:
        \begin{itemize}
            \item Roaming conflicts
            \item IP address conflicts
            \item Devices connecting to a weak/dead extender
        \end{itemize}
        \item \textbf{Action:} Locate and remove/replace this device
    \end{itemize}

    \item \textcolor{warningcolor}{\textbf{Channel 36 Congestion}}
    \begin{itemize}
        \item At least one competing network on same channel (-59 dBm)
        \item Can cause intermittent interference
        \item \textbf{Action:} Change router to channel 149, 153, or 157
    \end{itemize}
\end{enumerate}

\subsection{Medium Priority Issues}

\begin{enumerate}
    \item \textbf{4 Active VPN Tunnels (utun0-3)}
    \begin{itemize}
        \item Multiple tunnel interfaces competing for bandwidth
        \item May include iCloud Private Relay
        \item Can add latency and overhead
    \end{itemize}

    \item \textbf{MCS Index of 6 (Suboptimal)}
    \begin{itemize}
        \item WiFi 6 can achieve MCS 9-11 with better signal
        \item Indicates room for improvement in signal quality
    \end{itemize}
\end{enumerate}

\subsection{Low Priority (Currently Not Issues)}

\begin{itemize}
    \item VoIP devices: Currently idle, no impact
    \item Smart TVs: Idle, minimal bandwidth
    \item Printer: Minimal traffic
\end{itemize}

%==============================================================================
\section{Optimization Recommendations}
%==============================================================================

\subsection{Immediate Actions}

\begin{enumerate}
    \item \textbf{Remove/Replace Tenda Device}
    \begin{itemize}
        \item Locate the Tenda device at 192.168.50.11
        \item If it's a range extender, consider replacing with mesh WiFi
        \item Run: \texttt{sudo arp -d 192.168.50.11} to clear stale entry
    \end{itemize}

    \item \textbf{Change WiFi Channel}
    \begin{itemize}
        \item Access router at \url{http://192.168.50.1}
        \item Navigate to Wireless settings
        \item Change 5GHz channel from 36 to 149 or 157
    \end{itemize}

    \item \textbf{Update Router Firmware}
    \begin{itemize}
        \item ASUS regularly releases firmware updates
        \item Check for updates in router admin panel
    \end{itemize}
\end{enumerate}

\subsection{Recommended Settings}

\begin{center}
\begin{tabular}{lp{8cm}}
\toprule
\textbf{Setting} & \textbf{Recommendation} \\
\midrule
5GHz Channel & 149, 153, or 157 (less congested) \\
Channel Width & 80 MHz (current - good) \\
Band Steering & Enable (auto-select 5GHz when possible) \\
QoS & Enable for VoIP devices (prioritize .124, .150, .201) \\
DHCP Reservations & Set static IPs for all known devices \\
\bottomrule
\end{tabular}
\end{center}

\subsection{Network Architecture Improvements}

\begin{enumerate}
    \item \textbf{Replace Range Extender with Mesh}
    \begin{itemize}
        \item ASUS AiMesh compatible nodes work with RT-AX3000
        \item Eliminates roaming issues
        \item Seamless handoff between nodes
    \end{itemize}

    \item \textbf{Separate IoT Network (Optional)}
    \begin{itemize}
        \item Create guest network for IoT devices
        \item Isolates smart home devices from computers
        \item Reduces broadcast domain size
    \end{itemize}

    \item \textbf{VoIP QoS}
    \begin{itemize}
        \item Configure QoS to prioritize UDP ports 5060, 5061, 10000-20000
        \item Ensures call quality during high bandwidth usage
    \end{itemize}
\end{enumerate}

%==============================================================================
\section{Monitoring Commands}
%==============================================================================

Use these commands to monitor the network:

\begin{verbatim}
# Check router latency
ping 192.168.50.1

# Capture traffic conversations (30 seconds)
sudo tshark -i en0 -a duration:30 -qz conv,ip

# Protocol breakdown
sudo tshark -i en0 -a duration:20 -qz io,phs

# Check WiFi signal
/System/Library/PrivateFrameworks/Apple80211.framework/\
Versions/Current/Resources/airport -I

# Discover network devices
arp -a | grep 192.168.50

# Check for mDNS services
dns-sd -B _services._dns-sd._udp local.
\end{verbatim}

%==============================================================================
\section{Conclusion}
%==============================================================================

The \textbf{Superdooper 5G} network is generally healthy with good signal quality and no major issues detected during analysis. The primary concern is the \textbf{offline Tenda device} at 192.168.50.11, which should be investigated and likely removed.

\textbf{Summary of Actions:}
\begin{enumerate}
    \item Remove/replace Tenda range extender
    \item Change 5GHz channel to reduce congestion
    \item Consider mesh WiFi if coverage is needed
    \item Enable QoS for VoIP devices
\end{enumerate}

With these optimizations, network dropouts should be significantly reduced.

\end{document}
